% 原书第六篇6.4节第181--196题；印刷页253--254、PDF页262--263。
% 答案印刷347--348，PDF356--357；图6.10可用Debye积分数值替代低分辨率读数。
\moleculeproblem{181}{253}
考虑以下线形分子的全部可能独立振动模：
(a) $\mathrm{CO_2}$，O--C--O；
(b) $\mathrm{C_2H_2}$，H--C--C--H。
\exsolution
具有 $N$ 个原子的自由非线性分子有 $3N-6$
个内部振动自由度；线形分子只有两个独立转动轴，
故有 $3N-5$ 个。
(a) $\mathrm{CO_2}$ 有 $N=3$，
得到 $\boxed{3N-5=4}$：
两种沿分子轴的伸缩振动（对称与反对称）
及一种弯曲频率下两个互相垂直平面的简并弯曲模式。
(b) $\mathrm{C_2H_2}$ 有 $N=4$，
得到 $\boxed{3N-5=7}$：
三个沿分子轴的独立伸缩模式
及两种不同弯曲频率、各自具有两重横向简并的四个弯曲模式。
\excheck 两个横向弯曲自由度即使频率相等，
仍是\emph{两个}独立正常坐标，计数不可漏掉。

\moleculeproblem{182}{253}
长度为 $l$ 的理想细弦横波速度为 $v$，
所有横振动均处于同一平面。
求自然角频率区间 $(\omega,\omega+d\omega)$
内的正常振动数。
\exsolution
固定两端边界满足
$k_n=n\pi/l$，$\omega_n=vk_n$，
$n=1,2,\ldots$。
当模式密集时
\[
 N(\omega)\simeq\frac{l\omega}{\pi v},\qquad
 \boxed{dN=\frac l{\pi v}\,d\omega}.
\]
\excheck 角频率相邻间隔是 $\pi v/l$，
因此密度是其倒数；只有一个极化平面，不乘2。

\moleculeproblem{183}{253}
面积为 $S$ 的正方形薄膜，
垂直膜面的振动传播速度为 $v$。
求角频率区间 $(\omega,\omega+d\omega)$
内的固有振动数。
\exsolution
边长 $a=\sqrt S$，
本征波数
$k_{n_xn_y}=(\pi/a)\sqrt{n_x^2+n_y^2}$，
其中 $n_x,n_y>0$。
高频模式在 $k$ 空间的正象限内分布，
每格面积 $(\pi/a)^2$。
令 $k=\omega/v$，半径 $k$ 正象限圆面积
$\pi k^2/4$，故
\[
 N(\omega)\simeq
 \frac{S}{4\pi}\frac{\omega^2}{v^2},\quad
 \boxed{dN=\frac{S\omega}{2\pi v^2}\,d\omega}.
\]
\excheck 只计垂直膜面的一个位移极化分支，
不能擅自乘2。

\moleculeproblem{184}{253}
体积 $V$ 的直角平行六面体中，
横波传播速度 $v$。
求角频率区间 $(\omega,\omega+d\omega)$
内所有自然横振动的模式数。
\exsolution
高频下每一极化分支的驻波允许波数密度为
$V/(2\pi)^3$，三维球壳体积 $4\pi k^2dk$。
单分支为
\[
 dN_1=\frac{V}{2\pi^2}k^2dk
 =\frac{V\omega^2}{2\pi^2v^3}\,d\omega.
\]
横波有两种互相正交的独立极化，故
\[
 \boxed{dN=2dN_1=
 \frac{V\omega^2}{\pi^2v^3}\,d\omega}.
\]
\excheck 本题的两个横极化分支使答案恰为单分支的2倍。

\moleculeproblem{185}{253}
假定纵波、横波速度均为 $v$，
求理想各向同性 Debye 模型下
(a) 原子线密度 $n_0$ 的一维晶链；
(b) 原子面密度 $n_0$ 的二维正方晶格；
(c) 原子体积密度 $n_0$ 的三维简立方晶格
的 Debye 温度。
\exsolution
把 $D$ 维允许波数状态填至截止波数 $k_D$，
令所有声学分支的总模数等于原子数乘自由度数。
一维每原子一声学自由度，正、负波数满足
$k_D=\pi n_0$。
二维每原子两自由度，在波数圆盘内
$n_0=k_D^2/(4\pi)$，故 $k_D=\sqrt{4\pi n_0}$。
三维三声学分支满足
$n_0=k_D^3/(6\pi^2)$，故
$k_D=(6\pi^2n_0)^{1/3}$。
由 $\theta=\hbar vk_D/k_B$：
\[
 \boxed{\theta_{\rm1D}=\frac{\hbar v}{k_B}\pi n_0},\quad
 \boxed{\theta_{\rm2D}=\frac{\hbar v}{k_B}\sqrt{4\pi n_0}},\quad
 \boxed{\theta_{\rm3D}=\frac{\hbar v}{k_B}(6\pi^2n_0)^{1/3}}.
\]
\excheck 三维包含一纵两横声学支，
但每原子也有三个振动自由度，计数时两个3相消。

\moleculeproblem{186}{254}
铁中纵、横声速分别为
$v_l=5.85\,\mathrm{km/s}$ 与
$v_t=3.23\,\mathrm{km/s}$。
求其 Debye 温度。
\exsolution
三维各向同性晶体每单位体积声子态密度
\[
 \frac{dN}{d\omega\,V}=
 \frac{\omega^2}{2\pi^2}
 \left(v_l^{-3}+2v_t^{-3}\right).
\]
将截至 $\omega_D$ 的全部态数等于 $3n_0$，
得到
\[
 \omega_D^3=\frac{18\pi^2n_0}
 {v_l^{-3}+2v_t^{-3}},\qquad
 \boxed{\theta=\frac{\hbar}{k_B}
 \left(\frac{18\pi^2n_0}
 {v_l^{-3}+2v_t^{-3}}\right)^{1/3}}.
\]
取铁质量密度 $7.87\,\mathrm{g/cm^3}$、
摩尔质量 $55.85\,\mathrm{g/mol}$，
$n_0=\rho N_A/M$，求得
$\boxed{\theta\simeq4.71\times10^2\,\mathrm K}$。
\excheck 横波有两条极化分支；
不可对速度简单取算术平均值。

\moleculeproblem{187}{254}
铝的 Debye 温度 $\theta=396\,\mathrm K$。
忽略纵、横声速差异，
估算铝中声学振动传播速度 $v$。
\exsolution
等声速三维 Debye 截止满足
$\omega_D=v(6\pi^2n_0)^{1/3}$。
用 $\hbar\omega_D=k_B\theta$，
\[
 \boxed{v=\frac{k_B\theta}
 {\hbar(6\pi^2n_0)^{1/3}}}.
\]
铝的 $\rho\simeq2.70\,\mathrm{g/cm^3}$、
$M\simeq26.98\,\mathrm{g/mol}$，
$n_0=\rho N_A/M$，
得 $\boxed{v\simeq3.39\,\mathrm{km/s}}$。
\excheck 本题的结果是忽略极化差异的\emph{等效声速}，
不意味着真实纵波、横波速度相同。

\moleculeproblem{188}{254}
按一维 Debye 模型，
推导由相同原子组成的长链摩尔热容
$C_V(T)$，Debye 温度为 $\theta$；
并给出 $T\gg\theta$ 的极限。
\exsolution
一维声学模的角频率态密度在
$0<\omega<\omega_D$ 为常量，
每摩尔总模数 $N_A$，所以
$g(\omega)=N_A/\omega_D$。
每模平均能量为
$\hbar\omega/2+\hbar\omega/(e^{\hbar\omega/k_BT}-1)$，
因此
\[
 U=R\theta\left[\frac14+
 \left(\frac{T}{\theta}\right)^2
 \int_0^{\theta/T}\frac{x\,dx}{e^x-1}\right].
\]
令 $y=\theta/T$，求导得到
\[
 \boxed{C_V=R\left[
 \frac{2T}{\theta}\int_0^{\theta/T}
 \frac{x\,dx}{e^x-1}
 -\frac{\theta/T}{e^{\theta/T}-1}\right]}.
\]
高温下 $e^x-1\simeq x$，
故 $\boxed{C_V\longrightarrow R}$。
\excheck 一维晶链每原子只计一个纵向声学自由度，
因此高温极限为 $R$ 而非三维的 $3R$。

\moleculeproblem{189}{253--254}
长 $l$ 的单原子链晶格间距 $a$，
给定纵波色散关系
$\omega(k)=\omega_{\max}\sin(ka/2)$
（取第一布里渊区的非负频率支）。
求角频率态密度 $dN/d\omega$
及可取的纵振动总模数。
\exsolution
对周期边界的正、负传播波合计，
在非负波数变量 $k\in[0,\pi/a]$ 中，
态密度 $dN=(l/\pi)\,dk$。
从色散关系微分，
\[
 \frac{d\omega}{dk}=
 \frac{a\omega_{\max}}2
 \sqrt{1-\frac{\omega^2}{\omega_{\max}^2}}.
\]
因此
\[
 \boxed{\frac{dN}{d\omega}=
 \frac{2l}{\pi a\sqrt{\omega_{\max}^2-\omega^2}}}.
\]
积分 $0\le\omega\le\omega_{\max}$：
\[
 N=\frac{2l}{\pi a}
 \int_0^{\omega_{\max}}
 \frac{d\omega}{\sqrt{\omega_{\max}^2-\omega^2}}
 =\boxed{\frac la}.
\]
\excheck 带边群速度趋零，使态密度发散，
但该发散可积，总模数仍等于原子数。

\moleculeproblem{190}{254}
铜的 Debye 温度为 $\theta=330\,\mathrm K$。
求每克铜在 $T=0$ 时的零点振动能。
\exsolution
三维 Debye 模型一摩尔振动零点能
\[
 U_0=\int_0^{\omega_D}
 \frac{\hbar\omega}2
 \frac{9N_A\omega^2}{\omega_D^3}\,d\omega
 =\frac98N_A\hbar\omega_D
 =\frac98R\theta.
\]
按铜摩尔质量 $M=63.546\,\mathrm{g/mol}$，
每克能量
\[
 \boxed{\frac{U_0}{M}=
 \frac{9R\theta}{8M}
 \simeq48.6\,\mathrm{J/g}}.
\]
\excheck 零点能不因 $T\to0$ 消失，
其热容贡献却为零。

\moleculeproblem{191}{254}
原图6.10给出 Debye 晶体摩尔热容
$C/(3R)$ 随 $T/\theta$ 的曲线。
据此求：(a)银在 $T=65\,\mathrm K$、
$C=15\,\mathrm{J/(mol\,K)}$时的 $\theta$；
(b)铝在 $T=250\,\mathrm K$ 时
$C=22.4\,\mathrm{J/(mol\,K)}$，
求 $T=80\,\mathrm K$ 时热容；
(c)铜 $T=125\,\mathrm K$ 时热容较经典值低25\%，
求最大声子角频率。
\exsolution
无需从低分辨率原图人工读数，
可直接用原图所画的 Debye 函数：
\[
 \frac{C_V}{3R}=
 3\left(\frac{T}{\theta}\right)^3
 \int_0^{\theta/T}\frac{x^4e^x}{(e^x-1)^2}\,dx.
\]
(a) 设右式等于 $15/(3R)$，数值求根得
$\boxed{\theta_{\rm Ag}\simeq219\,\mathrm K\approx220\,\mathrm K}$。
(b) 从 $C_V(250)=22.4$ 求
$\theta_{\rm Al}\simeq371\,\mathrm K$，
再代 $80\,\mathrm K$，
$\boxed{C_V(80)\simeq10.3\,\mathrm{J/(mol\,K)}\approx10}$。
(c) 令 $C_V/(3R)=0.75$、
$T=125\,\mathrm K$，
数值求得 $\theta_{\rm Cu}\simeq309\,\mathrm K$。
故
\[
 \boxed{\omega_{\max}=\frac{k_B\theta}{\hbar}
 \simeq4.05\times10^{13}\,\mathrm{s^{-1}}}.
\]
\excheck (a)--(c)的差异仅为读图精度；
数值积分值与原图、书末约220K、10J/molK、$4.1\times10^{13}$
一致。

\moleculeproblem{192}{254}
从三维 Debye 模型的摩尔热容出发，
证明 $T\ll\theta$ 时
\[
 C_V=\frac{12\pi^4}{5}R(T/\theta)^3.
\]
\exsolution
由
\[
 C_V=9R(T/\theta)^3
 \int_0^{\theta/T}\frac{x^4e^x}{(e^x-1)^2}\,dx
\]
并令低温 $\theta/T\to\infty$，
利用分部积分
$\int_0^\infty x^4e^x/(e^x-1)^2dx
=4\int_0^\infty x^3/(e^x-1)dx$。
进一步展开
$(e^x-1)^{-1}=\sum_{j\ge1}e^{-jx}$，
\[
 4\int_0^\infty\frac{x^3}{e^x-1}dx
 =4\cdot3!\sum_{j=1}^\infty j^{-4}
 =24\zeta(4)=\frac{4\pi^4}{15}.
\]
所以
\[
 \boxed{C_V=\frac{12\pi^4}{5}
 R\left(\frac{T}{\theta}\right)^3}.
\]
\excheck 低温立方律是三维线性色散声子的态密度
$\propto\omega^2$ 的结果，而非普通经典热容极限。

\moleculeproblem{193}{254}
某晶体在 $20$ 与 $30\,\mathrm K$ 时
摩尔热容分别为 $0.226$ 与
$0.760\,\mathrm{J/(mol\,K)}$。
可否把这两个温度视为该晶体的 Debye 低温区？
\exsolution
低温 Debye 立方律要求
$C_V(T)\propto T^3$。
实测比例
\[
 \frac{C_V(30)}{C_V(20)}
 =\frac{0.760}{0.226}\simeq3.36,
\]
理论比例
$(30/20)^3=3.375$，非常接近。
因此
\[
 \boxed{\text{可以近似把20K和30K视为该晶体的Debye低温区}}.
\]
\excheck 此结论只检验热容是否近似 $T^3$，
不是对全部温度范围声子行为的判断。

\moleculeproblem{194}{254}
以 Debye 理论求零温下晶体\emph{每个振动振子}
的平均零点能，设 Debye 温度为 $\theta$。
\exsolution
三维 Debye 频率分布归一化为
每一个振动自由度
$g_1(\omega)=3\omega^2/\omega_D^3$。
每个频率为 $\omega$ 的振子零点能 $\hbar\omega/2$。
故
\[
 \overline E_0=
 \int_0^{\omega_D}\frac{\hbar\omega}{2}
 \frac{3\omega^2}{\omega_D^3}\,d\omega
 =\frac{3\hbar\omega_D}{8}
 =\boxed{\frac38k_B\theta}.
\]
\excheck 每个\emph{原子}有三声学振动自由度，
其平均零点能为 $9k_B\theta/8$；
本题问单个振子，要除以3。

\moleculeproblem{195}{254}
不计零点振动能，
画出 Debye 晶体在 $T=\theta/2$
与 $T=\theta/4$ 时单位频率区间的
振动能量分布曲线，并说明两曲线的比较。
\exsolution
三维 Debye 态密度
$g(\omega)\propto\omega^2$，
每个模热激发能
$\hbar\omega/(e^{\hbar\omega/(k_BT)}-1)$，
故
\[
 \boxed{\frac{dU_{\mathrm{th}}}{d\omega}
 =\frac{9N\hbar}{\omega_D^3}
 \frac{\omega^3}{e^{\hbar\omega/(k_BT)}-1}},
 \quad 0\le\omega\le\omega_D.
\]
令 $x=\omega/\omega_D$，
$T=\theta/2$ 时谱形
$f_2(x)=x^3/(e^{2x}-1)$；
$T=\theta/4$ 时谱形
$f_4(x)=x^3/(e^{4x}-1)$。
在 $0<x<1$ 的每一点有
$f_2(x)>f_4(x)$，
两曲线均从原点起始，随 $\omega$
在本有限频带内变化。
据两函数作图即可准确复现原书图47的
高温曲线在上、低温曲线在下的关系。
\excheck 此处已经给出完整可作图的解析函数；
若发表图稿，应据函数绘制坐标图，不用任意手画频谱。

\moleculeproblem{196}{254}
估算 Debye 温度 $\theta=330\,\mathrm K$ 的铜
中一个声子（题面译为声学量子）
可具有的最大能量与最大动量。
\exsolution
最大声子角频率
$\omega_D=k_B\theta/\hbar$，故
\[
 \boxed{E_{\max}=\hbar\omega_D=k_B\theta
 \simeq28.4\,\mathrm{meV}}.
\]
线性声学色散 $\omega=vk$，
声子准动量 $p=\hbar k=\hbar\omega/v$，
于是
\[
 \boxed{p_{\max}\simeq\frac{k_B\theta}{v}}.
\]
采用铜典型等效声速 $v\simeq3.5\,\mathrm{km/s}$
求得
$\boxed{p_{\max}\sim1.3\times10^{-24}\,
 \mathrm{kg\,m/s}
 =1.3\times10^{-19}\,\mathrm{g\,cm/s}}$。
\excheck 原书用词“photon (acoustic quantum)”易造成混淆；
这里是\emph{声子 phonon}而非光子 photon，
且没有给定声速，动量只能给量级并说明选用的$v$。
